This appendix holds two readings that the fitted models of \autoref{sec:four} do not provide, for two different reasons. The first is the work done on the measured photovoltaic series: what the record contains, how it is treated before anything is computed from it, and what is read off it. No model is fitted on that series at all, for the reason given in \autoref{sec:five}: it is the application domain of this study and not an experimental corpus. The second is the probe read by depth, which no model reports because Model~B absorbs the probe stage into a zero-mean offset by design. Nothing in this appendix is therefore a hypothesis test, and none of the decisions of \autoref{tab:bayesDecisions} depend on any part of it. The purpose is different and narrower: the synthetic design establishes what a patch grid can and cannot represent at frequencies fixed by $(f_s,P,S)$, and what follows asks what that arithmetic corresponds to in the one domain the report is written for, what the models do when handed the quantity the agent of \autoref{sec:nine} would be reading, and where inside the stack the arithmetic takes effect.

\subsection{What the record contains}

The source is a full calendar year of measurement at one-minute resolution, $525{,}601$ records timestamped in Central European Time and spanning the whole of 2017. Six channels are recorded: module power at a $30^\circ$ tilt, ambient temperature, global horizontal irradiance, global irradiance in the plane of array, wind speed and wind direction. The five weather channels are present for $99.1\%$ of the minutes and carry a large negative sentinel where an instrument returned nothing, which is mapped to a missing value on reading rather than treated as a measurement. Module power is the channel the study uses and the only one whose absence is structural: it is missing from $62.2\%$ of the records, overwhelmingly because the plant is not producing.

That last figure decides more than it first appears to. The design of \autoref{sec:six} reads a window of $544$ samples, a context of $480$ followed by a horizon of $64$, and a model cannot be handed a window interrupted by an absent sample. Taken literally, only $64$ of the $365$ days carry $544$ consecutive recorded power samples, and $97$ reach $480$. A day is therefore not a unit of data here but a container of fragments, and the longest unbroken fragment of each day is reported alongside it wherever a day is shown, because it is that number and not the day's total that determines whether the day can be used. Under the convention of the following subsection the constraint relaxes considerably, since a night is no longer a gap, but the distinction between a recorded value and a supplied one is retained throughout rather than dissolved into it.

\subsection{Missing samples and the zero convention}

Any use of this series has to say what a missing power sample means before it can say anything else. Classifying every missing sample by the solar elevation at its own timestamp separates two populations that are indistinguishable in the file itself. In $79.6\%$ of cases the sun is below the horizon: the plant is producing nothing, and zero is not an imputation but the measured quantity. A further $17.2\%$ fall in daylight gaps longer than fifteen minutes, of which seventy-eight runs exceed four hours and account for $49{,}204$ samples on their own; these are plant or logger outages during which nothing was produced and nothing was recorded, and zero stands for them as well. The elevation test is exact rather than heuristic at this site: on the fifteenth day of the year the sun is up between $8.10$ and $16.97$ hours and recorded samples run from $8.15$ to $17.10$, and the two boundaries agree to within a few minutes throughout the year.

The remaining $3.2\%$ are short dropouts inside a working production curve, $6{,}029$ runs with a median length of two minutes. A third of them are flanked by more than fifty watts and their ninetieth percentile flank is $166$ watts, so writing zero into them would stamp an edge down to zero and back within two minutes that the plant never produced. Two things make that edge worse than the absence it replaces. A rectangular notch is broadband in frequency, and a study of which frequencies a patch grid can and cannot represent should not manufacture content that lands on the candidate sites of \autoref{eq:locking-candidate}. More seriously, a zero written while irradiance remains valid and high is exactly the pattern \autoref{eq:inverter-trip} defines as \valueof{InverterTrip}, so the flat convention would lead the controller to isolate a photovoltaic side that never failed. These gaps are therefore filled by linear interpolation between their own recorded endpoints, which is both closer to the truth and smooth, and they are marked as supplied wherever the series is displayed.

The convention is an input convention for forecasting and nothing more. Supplying a value does not make it valid: $v_P(t)$ is unchanged by it, so a supplied sample still falsifies $V_t$, still asserts \valueof{SensorMalfunction} through \autoref{eq:sensor-malfunction}, and still routes the cycle to safe-mode selection. The forecaster receives a series in which absence reads as no production; the controller continues to receive the observability flag recording that the measurement was never made. Collapsing the two would replace an evidence-quality condition with a physical claim, which is the substitution \autoref{eq:sensor-malfunction} exists to prevent, and it is the reason the two representations of the same minute are kept apart rather than reconciled.

One further class of value is removed before any of this applies. Twenty samples across four days report between $1{,}000$ and $59{,}701$ watts against a $99.9$th percentile of $246$; one of them reads $43{,}730$ watts against a simultaneously measured irradiance of $32$ watts per square metre, which is not physically possible. They are recording glitches, and they matter beyond their count because the model standardises each context by its own mean and standard deviation, so a single such sample anywhere in a window flattens the production curve that window was meant to carry. They are treated as missing and then fall under the rules above.

\subsection{Reading a minute-resolution record in the units of the design}

The synthetic design is written at $f_s=512$\,Hz and this record is one sample per minute, and the two are reconciled without resampling, because the quantity carried between them is dimensionless. A tokeniser receives a sequence of samples and has no notion of seconds: a component whose period is $T$ samples, rather than $T$ seconds, completes $P/T$ cycles inside one patch and $S/T$ inside one stride, whatever the wall-clock rate those samples were taken at. Cycles per patch and cycles per stride are therefore the same quantities in both settings, and the lock condition, a whole number of cycles per patch or per stride, is the same predicate in both. A measured window is consequently supplied to a model unaltered and its sample index read as the design reads a sample index, and a component of period $T$ minutes appears at the cycles-per-patch value it would have had in the synthetic study. This is what makes the synthetic result transferable to this domain at all, and it is also what fixes which real periodicities are at risk: at the reference geometry a candidate site is a periodicity of sixteen minutes, of eight, of five minutes and twenty seconds, and so on down the harmonic series, which is the band in which cloud transients live rather than the band in which the daily cycle does.

\subsection{What is reported for a single day}

A day is shown as its recorded samples placed against time of day, each point coloured by its own value, so that the shape of the production curve and the dispersion around it are visible in the same picture: a thin arc is a clear day and a ragged cloud of points is a day of cloud transients. Each day also carries its longest unbroken fragment as a shaded span and a stated length, since that is the quantity that decides whether the day can be handed to a model. Several days are additionally drawn one under another on a single shared colour scale, which is the only form in which their magnitudes are comparable; a day whose transients reach above every other day's peak is visible there and is not visible in the per-day figures, each of which is normalised to itself.

\subsection{Forecasting one horizon ahead}

The forecast reading fixes a day and a time of day and asks what the next horizon holds. The horizon is the design's own, $64$ samples, which at this cadence is sixty-four minutes rather than an hour; the distinction is stated because the two are close enough to be conflated and the model's horizon is the fixed quantity. The context is the $2048$ samples immediately preceding the chosen time, taken by position and therefore crossing midnight backwards into the preceding days, so that a cut at four in the afternoon carries that day's production and the whole of the previous day's. Two reasons decide that length rather than a single calendar day: it is the context length the architecture accepts without truncation, and it is the window the retraining procedure of \autoref{sec:appRetraining} used, so the model is asked a question of the shape it was trained on. A multi-day context also contains more than one sunrise, which is the only arrangement in which the daily cycle is a component inside the window rather than the whole of it.

Two models are read on the identical context. The first is the retrained reference geometry of \autoref{tab:hfModels}, trained from scratch under the common budget on the synthetic mixture. The second is the published checkpoint at the same $(16,16)$ geometry, which is the stock configuration of the architecture and the reason $(16,16)$ is the reference geometry of the sweep. The comparison is included because it is the one the deployment question invites, and it is qualified here rather than left to be misread: the published checkpoint was pretrained on a large observational corpus that includes energy and weather series, the retrained one has seen only the synthetic mixture, and any difference between them on measured telemetry is in the first place a difference of training corpus. It is not evidence about structural aliasing in either direction, and it is reported on a single day at a single cut with no repetition and no interval.

\textcolor{red}{\textbf{Placeholder, to be replaced once the forecast figures are produced.} Comment here on the horizon figures, one paragraph per model and one for the pair. For each: whether the forecast tracks the measured decline over the sixty-four minutes or flattens against it, whether the median sits above or below the measurement and by how much relative to the day's peak, and whether the eighty-per-cent interval covers the measured values or is wider than the movement it is meant to bound. Then the pair: whether the two models differ mainly in level, in slope, or in the width of their intervals, and which of those differences the corpus argument of the preceding paragraph accounts for. State the horizon errors of both models with the number of measured points they were computed on, and state that they describe one cut on one day. Close by relating what the horizon shows to the periodicities named in the units subsection above: the transients the reference geometry places at candidate sites have periods of a few minutes, so a sixty-four-minute horizon contains several of them, and whether the forecast reproduces them at all is the observation this figure supports.}

\subsection{Frequency components of the measured series}

The forecast reading above is behavioural and says nothing about which components of the measured signal survive tokenisation. The decomposition reading asks that question directly, on the same series and in the same units, by comparing the frequency content of a measured window with the content of the window a model produces from it. It is the counterpart on real telemetry of the decomposition pages of \autoref{sec:appSweep}, which perform the same reading on generated signals whose composition is known in advance, and the two differ in exactly one respect that governs how each may be read: a synthetic window has a ground truth against which a recovered line is either correct or not, and a measured window does not, so what is available here is a comparison between two estimates rather than a measurement against a known value.

\textcolor{red}{\textbf{Placeholder, to be replaced once the decomposition on the measured series is complete.} This subsection reports what happens to the frequency content of a measured window when it passes through the tokeniser and the forecast. Describe first the input side: which lines the measured window carries once its mean is removed, where they fall in cycles per patch and cycles per stride at the reference geometry, and which of them coincide with the candidate sites $\mathcal F_{\mathrm{lock}}$ of that geometry within the resolution the window supports. State that resolution explicitly, since it is set by the window length and bounds every claim that follows, and state which part of the spectrum is the daily cycle and its harmonics, which are far below the candidate band and are not what this reading is about. Then the output side: which of those lines are present in the forecast, at what amplitude relative to the input, and whether any line appears that the input did not carry. Report the comparison at the reference geometry first and then across the geometries of \autoref{tab:hfModels}, so that a line lost at one geometry can be checked against a geometry whose combs fall elsewhere; this cross-geometry check is what distinguishes a component the tokenisation cannot represent from a component the plant simply stopped producing, and it is the only control available on measured data. Qualify the reading with the three limits that apply to it and do not apply to the synthetic decomposition: there is no ground-truth amplitude, so a recovered line is compared against an estimate carrying its own error; the measured content is not stationary across the window, since irradiance changes within any window long enough to resolve the candidate band, so a line may weaken for meteorological reasons rather than representational ones; and the supplied samples of the zero convention above contribute their own content, which is flat where a night was supplied and smooth where a short gap was interpolated, and must not be read as plant behaviour. Conclude by stating what the reading does and does not license: it is descriptive, it is not paired against controls and carries no interval, and the representational claim of \autoref{sec:four} is settled by the fitted estimands under the rule of \autoref{tab:bayesDecisions} and not here. Where the measured reading agrees with the synthetic result, say that it is consistent with it; where it does not, say which of the three limits above could account for the difference before treating it as a finding about the models.}

\subsection{Where the frequency stops being readable}

Model~B of \autoref{sec:appBayes} carries the probe stage as the zero-mean offset $s_{\text{st}}$, which absorbs the several-fold level differences between probe points so that they cannot be attributed to the lock indicator. That is the right treatment for the estimand and it leaves one question unanswered: the fit reports how much the labels cost at a candidate frequency, pooled over the stack, and says nothing about where inside the stack that cost appears. This subsection is the descriptive complement. It reports the same codelengths by probe stage rather than pooled, over the ten states named in \autoref{sec:appBayes}: the encoder $[\mathrm{REG}]$ token after each of the four encoder blocks, the decoder state after each of the four decoder blocks, the pre-projection vector and the output quantile head, in that order. No further measurement is required, because these are the cells Model~B is already fitted on.

The reading is descriptive in the same sense as the rest of this appendix. A depth profile is a picture of where a difference is largest, not a test that the difference exists; the representational claim of \autoref{sec:four} is settled by $e^{\theta_{\mathrm{lock}}}$ under the rule of \autoref{tab:bayesDecisions} and not by anything drawn here. The distinction matters more than usual in this case, because a depth profile is visually persuasive in a way a posterior interval is not.

\textcolor{red}{\textbf{Placeholder, to be replaced once the per-stage reading is produced.} Report the probe by depth in three paragraphs and two figures, plus an optional table. The first figure is the depth profile at the retrained reference geometry: the ten stages in order along the horizontal axis, mean codelength in bits on the vertical, one line for candidate frequencies and one for their controls, and a band showing the spread across frequencies within each stage. The quantity the figure exists for is the gap between the two lines as a function of depth, not the level of either. The paragraph that reads it should say three things whichever way the data fall: where the two lines separate, and whether that is early, in the encoder, so that the loss is at or near tokenisation, or late, in the decoder or the head, so that it is in the readout, or nowhere at all; whether the absolute level rises or falls with depth, which is a property of the representation size and not of the phenomenon; and that a gap confined to the output head would mean the projection rather than the encoder is where frequency identity is discarded, which is precisely the distinction a single-stage reading cannot make. If the two lines do not separate anywhere, that is a finding and the subsection is written the same way. The second figure repeats the profile across the geometries of \autoref{tab:hfModels}, and it should plot the locked-minus-control gap rather than the raw codelength, because Model~B blocks on geometry with $g_{\text{geo}}$ precisely for the reason that probe levels differ several-fold between configurations and are not comparable raw; state in the text that the figure leaves in the level differences the fit absorbs. The third paragraph reads that figure: whether the depth at which the gap opens is stable across $(P,S)$ or moves with the geometry, and then the two limits, stated here rather than left to be raised. The comparison is like-for-like in the probe task, the folds and the candidate frequencies; it is not like-for-like in the analysis context length, which already differs by geometry at $512$, $504$ and $500$ samples for the three stride groups of \autoref{sec:appSweep}, and saying so before a reader notices is worth more than the figure. The optional table has one row per stage and columns for the mean codelength at candidates, the mean at controls, the gap and the number of cells behind each, which makes the figures checkable at no cost. Finally, add no more than two sentences to \autoref{sec:seven}: one stating that the appendix reads the probe by depth so that the stage at which a candidate frequency stops being linearly separable can be located within the stack, and one stating that it is descriptive and that the representational claim is decided by $e^{\theta_{\mathrm{lock}}}$ under the rule of \autoref{tab:bayesDecisions}. The second sentence is not padding: it is what keeps the pointer from reading as a result.}
